\section{Introduction}

Let
\[
 H=\ker\!\left(\partial_x^2+\partial_y^2+\partial_z^2:
 \Q[x,y,z]_3\longrightarrow\Q[x,y,z]_1\right)
\]
be the rational seven-space of harmonic cubics for
\(Q(x,y,z)=x^2+y^2+z^2\).  We use the polarization pairing on cubics
normalized by
\[
 \langle p,(\,a\mathbin{\cdot}x\,)^3\rangle=p(a);
\]
its restriction to \(H\) is nondegenerate.  Hitchin associated to a
general \([f]\in\mathbf P(H)\) two icosahedral six-axis configurations:
equivalently, two Mukai--Umemura isotropic three-planes \(U\subset H\) with
\(f\perp U\) \cite{HitchinIcosahedron,HitchinBundles}.  The incidence
variety of pairs \(([f],U)\) is therefore generically a degree-two cover of
\(\mathbf P(H)\).  We determine its rational square class and, after fixing
the marked bridge datum below, compare its chosen sign with two further
realizations.  On the six-axis carrier the marked sign determines a
conference operator whose Pfaffian,
determinant, and middle-exterior shadows give the classical Joubert, Segre,
Igusa, and diagonal Clebsch models.  On the ten face axes it returns as the
degree-six Gaunt cubic.  Figure~\ref{fig:source-shadow-return} records this
architecture and the role of the marking.

\begin{figure}[H]
\centering
\resizebox{\linewidth}{!}{%
\begin{tikzpicture}[
  node distance=7mm and 13mm,
  box/.style={draw,rounded corners=2pt,align=center,inner sep=5pt,
    font=\small,text width=43mm},
  wide/.style={box,text width=55mm},
  arr/.style={-{Latex[length=2mm]},semithick},
  lab/.style={font=\scriptsize,fill=white,inner sep=1.5pt}
]
\node[box] (cover) {finite Stein normalization\\of Hitchin's cover\\[-1pt]
  \(z^2=5J_0\)};
\node[box,right=of cover] (datum) {marked bridge datum \(\mathfrak m\)\\
  chart lift, outer and Petersen labels};
\node[wide,below=15mm of $(cover)!0.5!(datum)$] (mark) {chosen golden sign\\
  with its compatible labels};
\node[box,below left=of mark] (six) {six ordered axes\\
  \(C^2=5I\), oriented cubic \(Z_{\mathfrak m}\)};
\node[box,below right=of mark] (ten) {ten labelled face axes\\
  Petersen four-space};
\node[box,below=of six] (classical) {four operator shadows\\
  Joubert--Segre--Igusa--Clebsch};
\node[box,below=of ten] (gaunt) {degree-six harmonic return\\
  \(\displaystyle \frac1{4\pi}\int F_y^3
     =-\frac{784000}{1247103}\sigma_3(y)\)};
\draw[arr] (cover) -- node[lab,above left] {pull back to \(\iota_t\)} (mark);
\draw[arr] (datum) -- node[lab,above right] {select component and label} (mark);
\draw[arr] (mark) -- (six);
\draw[arr] (mark) -- (ten);
\draw[arr] (six) -- (classical);
\draw[arr] (ten) -- (gaunt);
\end{tikzpicture}%
}
\caption{The source--shadow--return scheme.  Solid arrows are constructions
after the two top inputs have been fixed.  The lower branches compare the
chosen sign; they neither identify the displayed cubics nor define a map
between their ambient harmonic spaces.}
\label{fig:source-shadow-return}
\end{figure}

\paragraph{First-pass route.}
The main argument runs through the incidence cover in
Section~\ref{sec:cover}, its marked pullback in
Section~\ref{sec:orientation-source}, the golden fibre and exchanger in
Sections~\ref{sec:golden-fibre} and~\ref{sec:golden-exchanger}, the four operator shadows in
Section~\ref{sec:four-shadows}, and the harmonic return in
Section~\ref{sec:harmonic}.  The finite-field specialization in
Section~\ref{sec:finite-field-specialization} and the exchange-spectrum and
reconstruction subsections of Section~\ref{sec:golden-operator} are independent
consequences and may be skipped on a first reading.

Hitchin defines a real invariant sextic from the characteristic polynomial
of a quadratic moment operator.  His appendix rescales that operator while
computing its restriction, so we separate the geometric hypersurface from
its scale.  Let \(J_0\) denote the rational equation of Hitchin's irreducible
sextic hypersurface with the exact normalization
\[
 \iota_t^*J_0=16\sigma_3^2
\]
for the lifted Clebsch chart defined in Section~\ref{sec:cover}.  The
rationality and uniqueness of this normalization, and its relation to
Hitchin's real invariant, are checked there
\cite[opening theorem and Appendix]{HitchinIcosahedron}.
Write
\[
 V=\{(y_1,\ldots,y_5):\textstyle\sum_i y_i=0\},\qquad
 \sigma_3(y)=\frac13\sum_i y_i^3
\]
for the Clebsch four-space and its invariant cubic.
Here and below, \(D(\sigma_3)\) is the principal open set
\(\{\sigma_3\ne0\}\).
Because \(J_0\) is a section of \(\mathcal O_{\mathbf P(H)}(6)\), we use
\[
 K\bigl(\sqrt{cJ_0}\bigr),\qquad K=\Q(\mathbf P(H)),
\]
as shorthand: choose a nonzero rational section \(s\) of
\(\mathcal O_{\mathbf P(H)}(3)\), put \(j_s=J_0/s^2\in K^\times\), and take
\(K(\sqrt{cj_s})\).  Replacing \(s\) changes \(j_s\) by a square, so the
quadratic extension is unchanged.

\begin{theorem}[Arithmetic incidence cover]
\label{thm:arithmetic-main}
The function field of Hitchin's incidence cover is
\[
 \Q(\mathbf P(H))\bigl(\sqrt{5J_0}\bigr).
\]
If \(\mathcal N\to\mathbf P(H)\) is its finite Stein normalization, then,
after normalizing the anti-invariant summand, its algebra is
\[
 \mathcal O\oplus\mathcal O(-3),\qquad z^2=5J_0.
\]
For the Clebsch chart
\(\iota_t:\mathbf P(V_{\Q(\sqrt5)})\to\mathbf P(H_{\Q(\sqrt5)})\)
defined in Section~\ref{sec:cover},
\[
 \iota_t^*J_0=16\sigma_3^2.
\]
The chart is not defined over \(\Q\); Galois conjugation carries it to the
chart of the conjugate icosahedron.  The fibre over the rational point
\([xyz]\) is nevertheless a complete reduced fibre with residue algebra
\(\Q(\sqrt5)\).
\end{theorem}

The underlying Galois exchange of the two conjugate icosahedra is classical
\cite[Section~1.3]{HMSVOuter}.  The theorem identifies its occurrence here
with the square class of Hitchin's incidence cover.

The factor \(5\) has a concrete descent meaning: the cubic \([xyz]\) and
its unordered geometric fibre are rational, but choosing either
icosahedral configuration---equivalently, either conjugate Clebsch chart
through that point---requires \(\Q(\sqrt5)\).
This fibre is the model case for the descent: its residue algebra determines
the global quadratic twist in the proof of
Theorem~\ref{thm:arithmetic-main}.

The choice can be retained globally after normalizing the pullback to the
Clebsch chart, but its comparison with the conference and harmonic signs is
relative to a marking.  A \emph{marked bridge datum} \(\mathfrak m\) consists
of an ordering of the six golden projective axes, the five plane triples
labelled so that \(q_1=xyz\) and \(\sum q_i=0\), and the resulting normalized
linear lift
\[
 \iota_t(y)=\sum_i y_iq_i.
\]
The datum also identifies the ten face axes with the two-subsets of those
five labels.  Choose normalized representatives of the six projective axes.
They give a conference matrix
\(C_{\mathfrak m}\).  A different choice switches it by a diagonal sign
matrix.  The resulting switching-invariant triangle cubic is
\[
 Z_{\mathfrak m}(x)=\sum_{i<j<k}(C_{\mathfrak m})_{ij}
 (C_{\mathfrak m})_{jk}(C_{\mathfrak m})_{ki}x_ix_jx_k.
\]
It is translation-invariant and hence defines
\[
 Z_{\mathfrak m}\in\operatorname{Sym}^3
 \bigl((\Q^6/\Q\mathbf1)^*\bigr).
\]
Write \([C_{\mathfrak m},Z_{\mathfrak m}]_{\rm or}\) for this marked
source.  Changing the auxiliary representatives by \(a_i\mapsto d_i a_i\)
sends \(C_{\mathfrak m}\) to \(DC_{\mathfrak m}D\) and fixes
\(Z_{\mathfrak m}\).  The opposite source is represented by
\((-C_{\mathfrak m},-Z_{\mathfrak m})\).  The sign of the cubic generator
is retained, so the brackets are not projective notation.  Relabelling is
equivariant only when all labelled objects are transported together; it is
not quotiented silently.

\begin{proposition}[Relative marked orientation bridge]
\label{thm:orientation-source}
Fix a marked bridge datum \(\mathfrak m\).  Let
\(B=\mathbf P(V_E)\), where \(E=\Q(\sqrt5)\), and pull Hitchin's
incidence normalization back along \(\iota_t:B\to\mathbf P(H_E)\).
The normalization of this pullback is a disjoint union
\[
 B_+\amalg B_-.
\]
The fixed isotropic plane selects \(B_+\), and incidence deck exchange
interchanges the components.  For an odd generator normalized on \(B_+\),
their equations are
\[
 z=4\sqrt5\,\sigma_3,
 \qquad z=-4\sqrt5\,\sigma_3.
\]
On \(D(\sigma_3)\) they are the two distinct icosahedral configurations;
on the Clebsch cubic the unnormalized branches meet, while the normalization
remains disjoint.

Relative to \(\mathfrak m\), equip \(B_+\) with the source
\([C_{\mathfrak m},Z_{\mathfrak m}]_{\rm or}\) and equip \(B_-\) with its
deck-opposite.  At \([xyz]\), the golden
exchanger identifies this opposite with the conjugate configuration and
the representative
\((-C_{\mathfrak m},-Z_{\mathfrak m})\).  Under the marked primitive
pair-sum map
\[
 \beta:V\longrightarrow\Q^{\binom52},
 \qquad \beta(y)_{ij}=y_i+y_j,
\]
the same relative convention fixes the sign of the degree-six cubic in
Theorem~\ref{thm:harmonic-main}.  The normalized component alone is not
asserted to determine \(\mathfrak m\).
\end{proposition}

Here \(Z_{\mathfrak m}\) and \(\sigma_3\) are different cubics on spaces of
dimensions five and four.  The proposition compares their orientations relative
to \(\mathfrak m\); it does not identify their polynomial domains.

Theorem~\ref{thm:operator-shadows} is the operator stage on the main route.
The six outer translates of \(Z_{\mathfrak m}\) are simultaneously triangle
cubics, middle-exterior diagonals, commutator Pfaffians, and cross-golden
determinants.  Their Joubert coordinates land on the Segre cubic, whose
coordinate sections are diagonal Clebsch cubic surfaces; centered squaring is
the classical Segre--Igusa polar map
\cite[Sections~2.1--2.2]{HMSVOuter}.  The marked datum defines this operator
family.

\medskip
\noindent\textbf{Series perspective.}
\emph{Clebsch: Rigidity from Sparse Shadows} asks when incomplete invariants
determine the structures that produced them.  The present paper follows an
arithmetic incidence cover through its conference source and operator shadows
to a harmonic return.  Papers~I--III give three marked routes toward a common
cubic--operator package; Paper~V will prove their exact transports and reverse
recognition.  Paper~IV is the parallel lower branch: minimum-word pairs
reconstruct a marked conic plane and polarity, with no cubic identification
asserted.  Figure~\ref{fig:series-map} records this division.

\begin{figure}[ht]
\centering
\begin{tikzpicture}[
  x=1cm,y=1cm,>=Latex,
  paper/.style={draw,rounded corners=1.5pt,fill=black!2,
    align=center,font=\footnotesize,minimum height=8mm,text width=29mm,
    inner sep=2.5pt},
  active/.style={paper,very thick,fill=black!12},
  package/.style={paper,text width=38mm}
]
\node[paper]   (p1) at (0,1.45) {\textbf{I}\\deep-hole syndrome};
\node[paper]   (p2) at (0,0.35) {\textbf{II}\\quadratic trade};
\node[active]  (p3) at (0,-0.75) {\textbf{III}\\arithmetic/operator shadows};
\node[package] (gold) at (5.25,0.35) {marked cubic and\\golden operator};
\node[paper]   (p5) at (10.5,0.35) {\textbf{V}\\recognition and return};
\node[paper]   (p4) at (0,-2.05) {\textbf{IV}\\minimum-word pairs};
\node[package] (plane) at (5.25,-2.05) {marked conic plane\\and polarity};
\draw[dashed,->] (p1.east) -- ([yshift=3mm]gold.west);
\draw[dashed,->] (p2.east) -- (gold.west);
\draw[dashed,->] (p3.east) -- ([yshift=-3mm]gold.west);
\draw[<->,thick] (gold.east) -- (p5.west);
\draw[->,thick] (p4.east) -- (plane.west);
\end{tikzpicture}
\caption{Series map.  Dashed arrows mark the comparisons Paper~V will close;
its two-way solid arrow marks recognition and return.  The lower solid arrow
is Paper~IV's independent reconstruction.}
\label{fig:series-map}
\end{figure}

The proof follows the language changes in
Figure~\ref{fig:source-shadow-return}.  First, the rational incidence model and
its canonical classes give a quadratic field branched exactly along the
reduced sextic; the complete etale fibre at \([xyz]\) then determines its
square class \([5]\).  Next, Proposition~\ref{thm:orientation-source}
normalizes the chart pullback and, relative to the stated marking and lift,
follows deck exchange through the golden conference source and Petersen
pair-sum map.  Finally, Theorem~\ref{thm:operator-shadows} identifies the four
operator formulas, and Theorem~\ref{thm:harmonic-main} identifies the Petersen
coefficient module and computes the invariant cubic by one exact
spherical-moment evaluation.

Two later subsections of Section~\ref{sec:golden-operator} use the same
conference carrier but supply no hypothesis to this route.
Theorem~\ref{thm:balanced-exchange-rigidity} characterizes order six by
cut-independent balanced exchange spectrum.  Independently,
Theorem~\ref{thm:aligned-faithfulness} proves that aligned four-sets reconstruct
two-graphs from seven vertices onward, with seven sharp.  For conference
matrices the latter recovers the signing from order ten onward, up to switching
and global negation.

Hitchin proves the generic degree, identifies the branch sextic, constructs
the Clebsch chart, computes its restriction, and classifies the fibre over
\(xyz\)
\cite[\S\S2--4, 7--10]{HitchinIcosahedron}
\cite[\S\S4--5]{HitchinBundles}.  The arithmetic theorem fixes the remaining
rational constant in that square class and computes the spinor class of the
explicit specialization.  The Mukai--Umemura threefold supplies the
original threefold \cite{MukaiUmemura}; Hitchin supplies the incidence
construction and every formula used here.  Dye's finite theorem gives the
complementary square-\(5\) criterion for Clebsch hexagons
\cite{DyeClebsch}.  The abstract
equation \(z^2=5J_{\mathbf Z}\) is etale on \(D(10J_{\mathbf Z})\), but its
comparison with the geometric incidence variety is proved only after
inverting an unspecified integer.

The harmonic theorem places the classical degree-six bond-order
observable---the rotationally
invariant cubic formed from degree-six harmonic coefficients---
\cite{SNR1981,SNR1983} on the same four-space and computes its exact
restriction.  The coefficient factors as
\[
 -\frac{784000}{1247103}
 =-\frac{400}{46189}\frac{1960}{27}.
\]
The first factor is the universal \((6,6,6)\) Gaunt coupling
\(\bigl(\begin{smallmatrix}6&6&6\\0&0&0\end{smallmatrix}\bigr)^2\);
the second is the restriction scalar of the marked Petersen embedding.
Thus the prime \(11\) in the denominator belongs to the universal angular
momentum normalization, not to the golden descent.  The relative comparison
uses their four-dimensional \(A_5\)-modules and the invariant polynomial
\(\sigma_3\).

The decomposition of spherical harmonics under icosahedral symmetry is
classical, including the degree-six symmetry types used here
\cite{MeyerHarmonics,CohanHarmonics}; modern applications likewise treat
degree six as the lowest nonconstant even icosahedral channel
\cite{DharmavaramCapsids}.  Our calculation fixes a particular labelled
ten-face-axis zonal map, identifies its Petersen four-space, and evaluates
the cubic on that map.  It does not claim priority for the ambient
icosahedral decomposition or for the bond-order invariant.
